
\subsection{Reconcile a declared current-tax calculation}

The next finite calculation makes one further assumption explicit: a single
nonnegative rate applies to nonnegative taxable profit before credits or any
other regime-specific adjustments. This is a hypothetical arithmetic model,
named \texttt{uniform\_rate\_no\_credits\_v1}. It does not select a country's
tax rules. The author must select that model affirmatively; a false selection
is unsupported and an unknown selection remains unresolved.

Each declared account adjustment is now supported by identified allocation
items. An item taken from the completed return history must retain its original
event identity and signed amount, and must appear exactly once. Additional
items have explicit assumption references. The items assigned to an income
account must sum to that account's declared adjustment. Their category assignment
is supplied by the author; the arithmetic check does not establish its legal
appropriateness. In particular, the native return adjustment is not added again
after it has already been included in those account totals.

There are also four cross-field basis checks. At both opening and closing,
the declared asset bases of all lease rights must sum to the aggregate
right-of-use tax basis, and the declared liability bases must sum to the
aggregate lease-liability basis. The zero-valued right B participates in those
checks. A mismatch rejects the conditional calculation while preserving the
complete recorded financial report and the inconsistent declarations for
inspection. Neither carrying amount nor termination status supplies a tax basis.

\begin{samepage}
Let $P$ be recorded pretax income over all seven income accounts, $a_i$ the
identified signed adjustment items and $r$ the explicitly supplied rate. The
conditional taxable-profit expression and gross current tax are
\[
 T=P+\sum_i a_i,
 \qquad
 C=rT\quad\text{when }T\geq0.
\]
\end{samepage}
Here $C$ is gross conditional current tax under the selected arithmetic rule.
It is not an incremental expense, accrued payable or cash payment. The
calculation does not subtract opening accruals, apply credits, recognize
deferred tax or change the recorded closing state.

For the mixed example, $P=5$ and the retained return adjustment is $-24$.
Without any additional declared item, $T=5-24=-19$. The result is
\texttt{loss\_route\_required}; no negative tax amount or refund is computed.
A text declaration requesting a refund cannot change that result because it
does not supply an executable loss rule. Two hypothetical controls illustrate
the nonnegative branch. An expressly supplied additional adjustment of 19 gives
$T=0$ and $C=0$. An additional adjustment of 40 gives $T=21$ and, at the
illustrative rate $1/4$, $C=21/4$. Those extra adjustments are test assumptions,
not inferred company facts or recommended tax treatments.

The comparison retains the entire financial report, including cash of 361 and
recorded profit of 5, in every accepted arithmetic result. Inconsistent allocation
or basis declarations retain that same report alongside their rejection. Tests
also change native adjustment amounts, duplicate or omit identities, alter the
model finding and tamper with computed results. These checks establish the
specified arithmetic and preservation properties. A loss-utilization rule,
tax settlement, deferred recognition and annual close remain separate work.
